\clearpage
\setcounter{page}{1}
\maketitlesupplementary
\appendix

\section{Proofs}
\label{app:proofs}

Throughout, items $i\in[N+M]$ have fixed relevances $\tilde r_i\ge0$ and token values $v_i$, and $\mathcal B$ denotes
the token counts that the pruner supports, including zero.

\paragraph{Optimality of water-filling with discrete token counts.}
For a price $\lambda\ge0$, let $b_i^*(\lambda)\in\argmax_{b\in\mathcal B}\,\tilde r_iv_i(b)-\lambda b$ for every item,
and let $B^*=\sum_ib_i^*(\lambda)$ be the tokens that this allocation spends. We show that $\mathbf b^*(\lambda)$
maximizes $\sum_i\tilde r_iv_i(b_i)$ among all allocations with $b_i\in\mathcal B$ and $\sum_ib_i\le B^*$. Let
$\mathbf b$ be any such allocation. By the definition of $b_i^*(\lambda)$,
\begin{equation*}
  \tilde r_iv_i(b_i)-\lambda b_i\le\tilde r_iv_i\big(b_i^*(\lambda)\big)-\lambda b_i^*(\lambda)
\end{equation*}
for every item. Summing over the items gives
\begin{equation*}
  \sum_i\tilde r_iv_i(b_i)\le\sum_i\tilde r_iv_i\big(b_i^*(\lambda)\big)+\lambda\Big(\sum_ib_i-B^*\Big) ,
\end{equation*}
and the last term is not positive, since $\lambda\ge0$ and $\sum_ib_i\le B^*$. The argument uses neither the
concavity of $v_i$ nor real-valued token counts. The bisection of \cref{alg:iwf} returns the allocation of a price at
which at most $B$ tokens are spent, so this allocation is optimal for the tokens that it spends.

\paragraph{Proof of Proposition~\ref{prop:proportional}.}
Consider Eq.~\eqref{eq:relevance} with real-valued token counts and the token value $v(b)=b/(b+\tau)$, which is
increasing and strictly concave, with $v'(b)=\tau/(b+\tau)^2$. Since $v$ is increasing, a maximizer spends the whole
budget, and since the objective is concave and the constraint is linear, an allocation with $0<b_i<b_{\max}$ for
every frame is a maximizer if and only if there is a price $\lambda>0$ with $r_iv'(b_i)=\lambda$ for every frame.
Suppose that the proportional allocation $b_i=r_iB$ were a maximizer. By assumption it satisfies $0<b_i<b_{\max}$, so
\begin{equation*}
  \frac{r_i\,\tau}{(r_iB+\tau)^2}=\lambda\qquad\text{for every frame } i .
\end{equation*}
Every relevance $r_i$ would then be a root of the quadratic $\lambda(rB+\tau)^2-\tau r=0$ in $r$, which has at most
two distinct roots. This contradicts the assumption that $N\ge3$ relevances are pairwise distinct. Hence the
proportional allocation does not maximize Eq.~\eqref{eq:relevance}.

The same argument shows when a proportional allocation is optimal. The condition $r_iv'(r_iB)=\lambda$ holds for
every relevance only if $x\,v'(x)$ is constant, that is, only if $v(x)=a\log x+c$, a token value that never
saturates.

\paragraph{Proof of Proposition~\ref{prop:bound}.}
Let $c_i^{(t)}$ be the tokens of frame $f_i$ in the cache after turn $t$, with $c_i^{(0)}=0$. By construction,
$p_i^{(t)}=\max\big(0,b_i^{(t)}-c_i^{(t-1)}\big)$ and $c_i^{(t)}=\max\big(c_i^{(t-1)},b_i^{(t)}\big)$, so that
$p_i^{(t)}=c_i^{(t)}-c_i^{(t-1)}$ at every turn. The sum over turns telescopes,
\begin{equation*}
  \sum_{t=1}^{T}p_i^{(t)}=c_i^{(T)}-c_i^{(0)}=\max_{t}\,b_i^{(t)}\le b_{\max} ,
\end{equation*}
where the second equality unrolls the update of $c_i^{(t)}$. The prefill is well defined because the pruner ranks the
tokens of a frame once, independently of the query: the tokens that a later turn adds are exactly those ranked from
$c_i^{(t-1)}+1$ to $b_i^{(t)}$, and no cached token is ever replaced.

\section{Implementation Details}
\label{app:impl}

\paragraph{Candidates and relevance.} We implement \name and all baselines in PyTorch, and run all experiments on
two NVIDIA A100 GPUs. Every method draws its candidates from frames sampled at 1 fps and kept at their
native resolution. \name embeds frames with SigLIP~2 and transcript segments with BGE, and uses the question together
with its answer options as the query. The softmax temperature of the relevance, the diversity weight $\eta$, and the
text weight $\beta$ are tuned per benchmark.

\paragraph{Token value.} The measurements of \cref{fig:value} use 280 questions of LVBench~\cite{wang2024lvbench},
which is not among our evaluation benchmarks. For each question, the VLM answers from its eight most relevant frames
with every frame shown at the same token count, at nine levels from 4 tokens to the native token count.

\paragraph{Compute.} With $n$ prompt tokens, the prefill of a language model with $P$ parameters, $L$ layers, and
hidden size $d$ costs $2Pn+2n^2dL$ floating-point operations, where the first term covers the weights and the second
term covers attention. We sum this cost over every VLM call that a method makes for a question, including the
planning calls of AdaAlloc and the chain-of-thought and scorer calls of QuoTA. The vision encoder is not included; it
adds 48 TFLOPs per question for \name.

\paragraph{Baselines.} Top-K (w/ pruning) keeps 192 tokens for each of the top-64 frames, which fills the budget of
12,288 tokens. Q-Frame uses SigLIP~2 in place of its Long-CLIP scorer. For InternVL3.5-4B, whose frames have a
fixed cost of 256 tokens, the budget of 12,288 tokens corresponds to 48 frames.

\paragraph{Multi-turn evaluation.} We evaluate 200 conversations of MT-Video-Bench with 450 judged turns. The baselines
run as standard multi-turn chat, in which the KV cache of the whole past conversation is kept and the frames selected
for a new question are appended. The reuse statistics of \cref{fig:overlap}(b) replay the allocations of \name over
300 conversations at 12,288 tokens per turn.
